\documentclass[12pt,a4paper,oneside]{book}

\usepackage{cite}

\usepackage[utf8]{inputenc}
\usepackage[T1]{fontenc}
\usepackage[english]{babel}
\usepackage{amsmath}
\usepackage{amsfonts}
\usepackage{amssymb}
\usepackage{graphicx}
\usepackage{subfig}
\usepackage{fancyhdr}
\usepackage{appendix}
\usepackage{hyphenat}
\usepackage{pdfpages}

\usepackage{array,multirow,makecell}
\newcolumntype{C}[1]{>{\arraybackslash}p{#1}}

\usepackage{enumitem}
\setlist{leftmargin=*,itemsep=0pt}

\usepackage{centernot}
\usepackage[linesnumbered,ruled,vlined,english,onelanguage]{algorithm2e}

\usepackage{quotchap}
\makeatletter
\renewcommand{\@makechapterhead}[1]{
 \chapterheadstartvskip
 {\size@chapter{\sectfont\raggedright
 {\chapnumfont
 \ifnum \c@secnumdepth >\m@ne
 \if@mainmatter\thechapter
 \fi\fi
 \par\nobreak}
 {\raggedright\advance\leftmargin10em\interlinepenalty\@M #1\par}}
 \nobreak\chapterheadendvskip}}
\makeatother
\renewcommand*{\chapterheadendvskip}{\vspace{2cm}}

\usepackage{geometry}
\geometry{hmargin=2.5cm,vmargin=2.5cm}

\pagestyle{fancyplain}
\lhead{\fancyplain{}{\nouppercase{\textit{\leftmark}}}}
\chead{\fancyplain{}{}}
\rhead{\fancyplain{}{}}
\lfoot{\fancyplain{}{}}
\cfoot{\fancyplain{}{}}
\rfoot{\fancyplain{\thepage}{\thepage}}
\renewcommand{\headrulewidth}{1pt}
\renewcommand{\footrulewidth}{1pt}

\renewcommand{\thesection}{\arabic{section}}

\usepackage{titlesec}
\titleformat{\paragraph}{\fontsize{11}{10}\bfseries}{\theparagraph}{1em}{}
\titlespacing*{\paragraph}{0pt}{10pt plus 2pt minus 0pt}{0pt plus 2pt minus 0pt}

\setcounter{secnumdepth}{4}
\setcounter{tocdepth}{4}

\usepackage{array}
\usepackage{multirow}

\setlength{\parskip}{8pt}
\usepackage{setspace}

\usepackage{url}

\usepackage{hyperref}
\definecolor{darkblue}{rgb}{0.0, 0.0, 0.5}
\hypersetup{
 colorlinks,
 linktocpage=true,
 linkcolor={darkblue},
 citecolor={darkblue},
 urlcolor={blue}}

\sloppy

\author{PPP Team --- INSA University}
\title{Collaborative Threat Intelligence Platform with Blockchain}

\begin{document}
\pagenumbering{gobble}

% =============================================================
% COVER PAGE
% (Replace this with \includepdf[pages=-]{PageDeGarde.pdf} if you
%  have a pre-made institutional cover.)
% =============================================================
\begin{titlepage}
\centering
\vspace*{0.5cm}
{\large\textsc{INSA University}}\\[0.2cm]
{\large Ministry of Higher Education and Scientific Research}\\[1.5cm]

{\large\bfseries PERSONAL PROFESSIONAL PROJECT}\\[0.3cm]
{\normalsize submitted in partial fulfilment of the\\
3\textsuperscript{rd} Year requirements in Computer Engineering}\\[1.8cm]

\rule{\linewidth}{0.6pt}\\[0.4cm]
{\Huge\bfseries A Collaborative Threat Intelligence Platform\par}
\vspace{0.4cm}
{\LARGE\itshape Enhancing Trust and Integrity with Blockchain Technology\par}
\rule{\linewidth}{0.6pt}\\[2cm]

{\large\textbf{Prepared by the PPP Team}}\\[0.2cm]
{\normalsize Group of 4 students --- Computer Engineering}\\[1.5cm]

{\normalsize Defended in front of a jury composed of:}\\[0.4cm]
\begin{tabular}{ll}
Jury President & Prof. [Name] \\
Reviewer       & Prof. [Name] \\
Supervisor at INSA & Prof. [Name] \\
\end{tabular}

\vfill
{\large Academic Year 2025--2026}
\end{titlepage}

% =============================================================
% ACKNOWLEDGEMENTS
% =============================================================
\chapter*{Acknowledgements}
\markboth{Acknowledgements}{}

We would like to express our sincere gratitude to all the people who supported
us throughout this Personal Professional Project. Our deepest thanks go to
the academic staff of INSA University, and in particular to the teachers of
the Computer Engineering program, whose guidance, technical insight and
constant encouragement allowed us to carry this project from an initial idea
to a working platform.

We are particularly grateful to our project supervisor for the time invested
in reviewing our work, for the constructive feedback offered during each
iteration and for the autonomy granted to us in technical decision-making.
This trust was a decisive factor in the depth we were able to reach.

We also wish to thank the members of the cybersecurity and open-source
communities whose publications --- on threat sharing platforms, blockchain
registries and large language models --- shaped the technical orientation
of this project.

Finally, our gratitude goes to our families and friends for their unwavering
support, and to one another, for the spirit of collaboration that made this
team work both productive and rewarding.

\frontmatter

% =============================================================
% ABSTRACT
% =============================================================
\chapter*{Abstract}
\markboth{Abstract}{}

\normalsize{This report presents the design and implementation of a
collaborative threat intelligence platform enhanced with blockchain
technology. Faced with the rapid growth of cyber threats and the difficulty
for organisations to share indicators of compromise in a trustworthy manner,
this project proposes a centralised solution that aggregates, validates and
exposes technical indicators, threat actor profiles and malware samples. The
platform is built on a layered architecture combining a FastAPI backend, a
PostgreSQL database, a conversational assistant powered by a large language
model, and a blockchain layer responsible for the immutable anchoring of
validated indicators. The report describes the project context, the
theoretical background of the involved technologies, the design and
implementation of the prototype, and the evaluation of the obtained results
against the initial objectives.

\medskip
{\noindent \textbf{Keywords: } cybersecurity, threat intelligence,
blockchain, indicators of compromise, smart contract, large language model,
FastAPI, Solidity. }

% =============================================================
% RÉSUMÉ
% =============================================================
\chapter*{Résumé}
\markboth{Résumé}{}

\normalsize{Ce rapport présente la conception et la réalisation d'une
plateforme collaborative de \textit{Threat Intelligence} enrichie par la
technologie \textit{blockchain}. Face à la prolifération des cybermenaces et
à la difficulté pour les organisations de partager leurs indicateurs de
compromission de manière fiable, ce projet propose une solution centralisée
permettant l'agrégation, la validation et la consultation d'indicateurs
techniques, de profils d'acteurs malveillants et d'échantillons de logiciels
malveillants. La plateforme repose sur une architecture en couches articulant
un \textit{backend} FastAPI, une base de données PostgreSQL, un assistant
conversationnel basé sur un modèle de langage et une couche
\textit{blockchain} dédiée à l'ancrage immuable des indicateurs validés.

\medskip
{\noindent \textbf{Mots-clés : } cybersécurité, \textit{threat intelligence},
\textit{blockchain}, indicateurs de compromission, \textit{smart contract},
modèle de langage, FastAPI, Solidity. }

\spacing{1}
\tableofcontents{}
\newpage
\listoffigures
\newpage
\listoftables
\newpage
\spacing{1.4}

% =============================================================
% LIST OF ACRONYMS
% =============================================================
\chapter*{List of Acronyms}
\markboth{List of Acronyms}{}
\begin{itemize}
\item \textbf{AI} Artificial Intelligence
\item \textbf{API} Application Programming Interface
\item \textbf{APT} Advanced Persistent Threat
\item \textbf{C2} Command and Control
\item \textbf{CTI} Cyber Threat Intelligence
\item \textbf{DLT} Distributed Ledger Technology
\item \textbf{EVM} Ethereum Virtual Machine
\item \textbf{HTTP} Hypertext Transfer Protocol
\item \textbf{IOC} Indicator of Compromise
\item \textbf{JSON} JavaScript Object Notation
\item \textbf{JWT} JSON Web Token
\item \textbf{LLM} Large Language Model
\item \textbf{MISP} Malware Information Sharing Platform
\item \textbf{NLP} Natural Language Processing
\item \textbf{ORM} Object-Relational Mapping
\item \textbf{REST} Representational State Transfer
\item \textbf{SDK} Software Development Kit
\item \textbf{SHA} Secure Hash Algorithm
\item \textbf{SIEM} Security Information and Event Management
\item \textbf{SQL} Structured Query Language
\item \textbf{TLP} Traffic Light Protocol
\item \textbf{TTP} Tactics, Techniques and Procedures
\item \textbf{UI} User Interface
\item \textbf{UUID} Universally Unique Identifier
\end{itemize}

\mainmatter

% =============================================================
% GENERAL INTRODUCTION
% =============================================================
\chapter*{General Introduction}
\addcontentsline{toc}{chapter}{General Introduction}
\markboth{General Introduction}{}

The digital transformation of public administrations, private companies and
critical infrastructures has multiplied the surface of exposure to cyber
threats. Attackers no longer act in isolation: malicious campaigns are
organised, automated and propagated across organisational boundaries within
minutes. A piece of attacker infrastructure --- a command-and-control server,
a phishing domain or a malware signature --- detected by one organisation
today is very likely to target another organisation tomorrow. Cyber defence
has therefore become a fundamentally collaborative discipline.

This collaborative dimension materialises around the concept of Cyber Threat
Intelligence (CTI). At the heart of CTI lie the Indicators of Compromise
(IOCs), atomic pieces of evidence that can be used by detection tools to
identify malicious activity. Yet, while large sharing communities exist, the
act of trusting an external indicator remains a challenge: how can an analyst
be sure that a published indicator has not been tampered with after
publication? How can a contributing organisation prove the authenticity of
what it submitted? These two questions --- the trustworthiness of data and
the immutability of the historical record --- are at the core of the project
presented in this report.

Within the framework of our Personal Professional Project at INSA University,
we designed and developed an academic platform that combines three
complementary ingredients. First, a classical web application providing
structured submission, storage and search of threat intelligence artefacts.
Second, a conversational assistant based on a large language model, intended
to lower the barrier of entry for non-expert visitors. Third, an integration
with a blockchain network that anchors every validated indicator to a
tamper-evident ledger.

The remainder of this report is organised into three balanced chapters. The
first chapter establishes the project scope: it describes the academic
context, formulates the problem and derives the functional and non-functional
requirements of the platform. The second chapter is dedicated to the
theoretical background of the three technological pillars that compose the
solution: collaborative threat intelligence, blockchain technology and
large language models. The third and final chapter presents the design,
the implementation choices and a structured evaluation of the delivered
prototype. A general conclusion summarises the contributions and outlines
the perspectives identified for future work.

% =============================================================
% CHAPTER 1 — PROJECT SCOPE
% =============================================================
\chapter{Project Scope}
\label{ch:scope}
\section*{Introduction}
This opening chapter delimits the scope of the project. It presents the
academic context within which the work was conducted, analyses the problem
addressed by the platform, surveys the existing solutions on the market and
finally derives the set of functional and non-functional requirements that
will guide the rest of the report.

% ---------------- Section 1
\section{Project Context}

\subsection{Academic Framework}

The work presented in this report was carried out in the context of the
Personal Professional Project (PPP), a structured exercise placed at the
heart of the third year of the Computer Engineering program at INSA
University. The PPP is intended to confront students, organised in small
teams, with a complete software engineering experience: identifying a real
problem, designing a system, choosing the right technologies, implementing
the solution and finally evaluating the result. Unlike a classroom assignment,
the PPP is sized to span an entire semester and explicitly rewards autonomy,
methodology and the ability to integrate heterogeneous components.

\subsection{Team Organisation}

The team is composed of four third-year students. The work was organised
into five sequential phases, each placed under the responsibility of a
dedicated lead while requiring contributions from every member. Weekly
synchronisation meetings, a shared Git repository and progressive integration
between modules allowed the team to converge despite the diversity of the
topics handled in parallel.

\subsection{Project Schedule}

The development effort is articulated around the 14-week schedule illustrated
in Figure \ref{fig:gantt}. Each phase ends with an internal demonstration
that conditions the start of the next phase.

\begin{figure}[!h]
\centering
\includegraphics[width=0.85\textwidth]{images/gantt.png}
\caption{Gantt chart of the development phases over a 14-week semester. Source: project planning document.}
\label{fig:gantt}
\end{figure}

The five phases and their main deliverables are summarised in
Table \ref{tab:phases}.

\begin{table}[!h]
\caption{Development phases and main deliverables.}
\label{tab:phases}
\centering
\renewcommand{\arraystretch}{1.4}
\begin{tabular}{| C{.08\textwidth} | C{.12\textwidth} | p{.65\textwidth} |}
\hline
\textbf{Phase} & \textbf{Weeks} & \textbf{Main deliverables} \\
\hline\hline
1 & S1--S3   & \nohyphens{Backend skeleton, database schema, API key authentication} \\
\hline
2 & S4--S7   & \nohyphens{Conversational assistant and TLP filtering middleware} \\
\hline
3 & S8--S9   & \nohyphens{Public dashboard and administration interface} \\
\hline
4 & S10--S11 & \nohyphens{Blockchain learning and smart contract prototyping} \\
\hline
5 & S12--S14 & \nohyphens{Backend integration with the ledger and final tests} \\
\hline
\end{tabular}
\end{table}

\subsection{Working Methodology}

The team followed an incremental approach inspired by agile principles. A
short Kanban board, hosted on a shared workspace, materialised the backlog
and made the workload of each member visible. Code reviews systematically
preceded every merge into the main branch, in order to enforce a homogeneous
coding style and to share technical knowledge across the team.

% ---------------- Section 2
\section{Problem Analysis}

\subsection{The Evolving Cybersecurity Landscape}

Cybersecurity is no longer a peripheral concern confined to a few
specialised teams. Recent reports from international agencies show a steady
increase in the number, sophistication and impact of incidents. Ransomware
campaigns disrupt hospitals, phishing operations target banking customers at
scale, and state-sponsored groups penetrate strategic supply chains for
long-term espionage. The cost associated with data breaches has reached
unprecedented levels, and the regulatory pressure --- through frameworks
such as the GDPR or the NIS2 directive in Europe --- forces organisations to
adopt a proactive posture.

\subsection{Problem Statement}

In this hostile environment, no single organisation can defend itself in
isolation. Information about attacker infrastructure must circulate quickly
between the defenders. However, current information-sharing communities
suffer from three recurring weaknesses that motivated this project.

\begin{itemize}
\item \textbf{Absence of cryptographic proof of integrity.} Once an
indicator is published on a sharing platform, nothing prevents an internal
modification of the historical record. A contributor cannot prove that what
was published yesterday is exactly what is being displayed today.
\item \textbf{High barrier of entry for non-experts.} The interfaces of
established platforms are intentionally dense and require knowledge of the
underlying data model, which limits adoption by smaller organisations.
\item \textbf{Heterogeneous trust models.} Contributors join through
bilateral agreements or through opaque approval processes, which makes it
difficult for a newcomer to know how much weight to give to a published
indicator.
\end{itemize}

These observations crystallise into the central question of this project:

\textbf{How can we design and implement a collaborative threat intelligence
platform that is open to public consultation, secure for contributors,
accessible to non-experts thanks to an AI assistant, and that provides a
verifiable cryptographic proof of integrity for every validated indicator?}

\subsection{Survey of Existing Solutions}

Several mature platforms already address the problem of threat intelligence
sharing. The most widely deployed are summarised in Table
\ref{tab:existing}.

\begin{table}[!h]
\caption{Comparison of existing collaborative threat intelligence platforms. Source: official project documentation and academic surveys.}
\label{tab:existing}
\centering
\renewcommand{\arraystretch}{1.4}
\begin{tabular}{| C{.16\textwidth} | p{.50\textwidth} | C{.10\textwidth} | C{.12\textwidth} |}
\hline
\textbf{Platform} & \textbf{Main purpose} & \textbf{Open source} & \textbf{Integrity proof} \\
\hline\hline
MISP & \nohyphens{Structured sharing of IOCs between trusted partners} & Yes & No \\
\hline
OpenCTI & \nohyphens{Aggregation and visualisation of CTI knowledge graphs} & Yes & No \\
\hline
AbuseIPDB & \nohyphens{Public crowd-sourced IP reputation database} & No & No \\
\hline
VirusTotal & \nohyphens{File and URL scanning, multi-engine analysis} & No & No \\
\hline
\end{tabular}
\end{table}

These platforms confirm the relevance of the collaborative model, yet none
of them addresses simultaneously the three weaknesses listed above. This gap
defines the contribution targeted by the present project.

\subsection{Proposed Solution}

To answer the problem statement, we propose an academic platform organised
around three pillars: a structured submission workflow accessible through
an HTTP API, a conversational AI assistant designed to lower the barrier
of entry, and a blockchain anchoring layer dedicated to integrity. The
detailed design is the subject of Chapter \ref{ch:solution}.

% ---------------- Section 3
\section{Solution Requirements}

\subsection{General Objectives}

The objectives derived from the problem statement are summarised in
Table \ref{tab:objectives}. They will be revisited in
Chapter \ref{ch:solution} when the prototype is evaluated.

\begin{table}[!h]
\caption{General objectives of the project.}
\label{tab:objectives}
\centering
\renewcommand{\arraystretch}{1.4}
\begin{tabular}{| C{.20\textwidth} | p{.70\textwidth} |}
\hline
\textbf{Objective} & \textbf{Description} \\
\hline\hline
Centralisation & \nohyphens{Aggregate IOCs, threat actors and malware samples submitted by multiple organisations into a single repository.} \\
\hline
Accessibility & \nohyphens{Allow any visitor to consult public-level intelligence without creating an account.} \\
\hline
AI assistance & \nohyphens{Provide a chatbot able to answer basic questions about cybersecurity concepts and to guide users through the database.} \\
\hline
Integrity & \nohyphens{Anchor each validated indicator on a blockchain ledger, providing a verifiable proof of non-tampering.} \\
\hline
Security & \nohyphens{Authenticate contributing organisations exclusively through API keys, with traceable usage and instant revocation.} \\
\hline
Compliance & \nohyphens{Enforce the Traffic Light Protocol to control the visibility of submitted data.} \\
\hline
\end{tabular}
\end{table}

\subsection{Actors and Roles}

The platform serves three categories of users with clearly differentiated
permissions, illustrated in Figure \ref{fig:actors}.

\begin{figure}[!h]
\centering
\includegraphics[width=0.85\textwidth]{images/usecase.png}
\caption{Simplified use-case diagram of the three categories of users. Source: authors.}
\label{fig:actors}
\end{figure}

\subsubsection{Public Visitor}
Any unauthenticated user can browse the public dashboard, search indicators,
consult their details and interact with the conversational assistant. The
visitor never sees data classified above the TLP:GREEN/WHITE levels.

\subsubsection{Contributor Organisation}
A contributor is a verified organisation that submits intelligence data
through a dedicated API key. The platform deliberately avoids a classical
login-password model: each request is identifiable, traceable and a
compromised key can be revoked instantly.

\subsubsection{Administrator}
The administrator is the only role that uses session-based authentication
(e-mail, password and a JSON Web Token). The administrator approves new
contributors, generates and revokes API keys, and supervises pending
submissions.

\subsection{Functional Requirements}

The functional requirements derived from the use-case analysis are
summarised in Table \ref{tab:fr}.

\begin{table}[!h]
\caption{Main functional requirements of the platform.}
\label{tab:fr}
\centering
\renewcommand{\arraystretch}{1.35}
\begin{tabular}{| C{.08\textwidth} | p{.65\textwidth} | C{.17\textwidth} |}
\hline
\textbf{Code} & \textbf{Requirement} & \textbf{Actor} \\
\hline\hline
FR-1  & \nohyphens{Submit a new IOC together with its descriptive context} & Contributor \\
\hline
FR-2  & \nohyphens{Submit a threat actor profile} & Contributor \\
\hline
FR-3  & \nohyphens{Submit a malware sample} & Contributor \\
\hline
FR-4  & \nohyphens{Browse and search the public catalogue of indicators} & Visitor \\
\hline
FR-5  & \nohyphens{Consult the detail and enrichment of an indicator} & Visitor \\
\hline
FR-6  & \nohyphens{Interact with the AI assistant in natural language} & Visitor / Contrib. \\
\hline
FR-7  & \nohyphens{Verify the blockchain proof of an indicator} & Visitor \\
\hline
FR-8  & \nohyphens{Flag one's own IOC as a false positive} & Contributor \\
\hline
FR-9  & \nohyphens{Approve, refuse or revoke a contributor application} & Administrator \\
\hline
FR-10 & \nohyphens{Review the submissions placed in a pending state} & Administrator \\
\hline
\end{tabular}
\end{table}

\subsection{Non-Functional Requirements}

The quality attributes that act as guidelines during the design phase are
reported in Table \ref{tab:nfr}.

\begin{table}[!h]
\caption{Non-functional requirements considered for the project.}
\label{tab:nfr}
\centering
\renewcommand{\arraystretch}{1.4}
\begin{tabular}{| C{.20\textwidth} | p{.70\textwidth} |}
\hline
\textbf{Attribute} & \textbf{Expected level} \\
\hline\hline
Security & \nohyphens{All submissions are authenticated by API key; keys are stored hashed and never in clear; TLP filters are enforced server-side.} \\
\hline
Integrity & \nohyphens{Every validated IOC is anchored on a blockchain; any modification is detectable.} \\
\hline
Performance & \nohyphens{Public endpoints answer in less than one second under nominal load.} \\
\hline
Scalability & \nohyphens{The architecture supports the addition of new contributors without redesign.} \\
\hline
Maintainability & \nohyphens{The code base is modular, documented and supported by automated database migrations.} \\
\hline
Auditability & \nohyphens{Every contributor request is logged with a timestamp, source organisation and operation type.} \\
\hline
Usability & \nohyphens{The public interface is understandable by a non-technical visitor in less than three minutes.} \\
\hline
\end{tabular}
\end{table}

\section*{Conclusion}
This first chapter established the scope of the project. It described the
academic context in which the work was conducted, analysed the problem of
collaborative threat intelligence and surveyed the existing platforms. From
this analysis, a precise problem statement was formulated, and the
functional and non-functional requirements that will drive the rest of the
report were enumerated. The next chapter dives into the theoretical
background of the three technological pillars on which the proposed
solution is built.

% =============================================================
% CHAPTER 2 — SOLUTION BACKGROUND
% =============================================================
\chapter{Solution Background}
\label{ch:background}
\section*{Introduction}
Before describing how the platform was designed and built, it is necessary
to introduce the conceptual foundations on which the work relies. This
chapter is dedicated to the theoretical background of the three pillars of
the solution. The first section addresses the discipline of cyber threat
intelligence, the second section explores blockchain technology and smart
contracts, and the third section presents the role of large language models
in conversational assistance. Each section progresses from the general
definition of the concept to the specific design choices that are relevant
to the project.

% ---------------- Section 1
\section{Cyber Threat Intelligence}

\subsection{Definition and Levels of Intelligence}

Cyber Threat Intelligence (CTI) can be defined as the collection and
analysis of structured information regarding adversaries and their
capabilities. The aim is to convert raw data into actionable knowledge that
supports better decisions. The literature traditionally distinguishes three
levels of intelligence, summarised in Table \ref{tab:cti-levels}.

\begin{table}[!h]
\caption{The three classical levels of cyber threat intelligence.}
\label{tab:cti-levels}
\centering
\renewcommand{\arraystretch}{1.4}
\begin{tabular}{| C{.18\textwidth} | C{.18\textwidth} | p{.50\textwidth} |}
\hline
\textbf{Level} & \textbf{Audience} & \textbf{Content} \\
\hline\hline
Strategic    & Executives        & \nohyphens{Long-term trends, geopolitical context, risk evolution.} \\
\hline
Tactical     & Security managers & \nohyphens{Adversary tactics, techniques and procedures (TTPs).} \\
\hline
Operational  & SOC analysts      & \nohyphens{Concrete indicators ready to be ingested by detection tools.} \\
\hline
\end{tabular}
\end{table}

The project focuses on the operational level, which is the one most directly
served by collaborative platforms.

\subsection{Indicators of Compromise}

At the operational level, the central artefact is the Indicator of
Compromise (IOC): a piece of technical evidence whose presence on an
information system is considered a strong sign of malicious activity. Table
\ref{tab:ioc-types} lists the main categories of IOCs handled in this
project.

\begin{table}[!h]
\caption{Main categories of indicators of compromise considered in the project.}
\label{tab:ioc-types}
\centering
\renewcommand{\arraystretch}{1.4}
\begin{tabular}{| C{.18\textwidth} | p{.72\textwidth} |}
\hline
\textbf{Type} & \textbf{Description and example} \\
\hline\hline
IP address & \nohyphens{Network endpoint involved in malicious activity (scanning, brute-force, command-and-control). Example: 185.220.101.45.} \\
\hline
URL & \nohyphens{Web resource hosting phishing pages, payloads or exfiltration endpoints.} \\
\hline
File hash & \nohyphens{Cryptographic fingerprint (MD5, SHA-256) of a known malicious binary.} \\
\hline
E-mail address & \nohyphens{Sender address used in phishing or business e-mail compromise campaigns.} \\
\hline
\end{tabular}
\end{table}

Individually, an IOC has a limited lifespan because attackers rotate their
infrastructure. Collectively, however, large feeds of indicators become a
strategic asset that can be integrated into firewalls, intrusion detection
systems, SIEM platforms and endpoint protection solutions.

\subsection{The Traffic Light Protocol}

Sharing intelligence between organisations raises the question of
confidentiality: not all data can circulate as openly as a public weather
forecast. The Traffic Light Protocol (TLP) is an international standard
defined by FIRST.org for marking sensitive information. Its four levels are
recalled in Table \ref{tab:tlp}.

\begin{table}[!h]
\caption{Confidentiality levels defined by the Traffic Light Protocol.}
\label{tab:tlp}
\centering
\renewcommand{\arraystretch}{1.4}
\begin{tabular}{| C{.15\textwidth} | p{.55\textwidth} | C{.15\textwidth} |}
\hline
\textbf{Level} & \textbf{Meaning} & \textbf{Public dashboard} \\
\hline\hline
TLP:RED   & \nohyphens{Strictly restricted to named recipients} & No \\
\hline
TLP:AMBER & \nohyphens{Restricted to the organisation and its partners} & No \\
\hline
TLP:GREEN & \nohyphens{Shared inside the cybersecurity community} & Yes \\
\hline
TLP:WHITE & \nohyphens{Public information, no restriction} & Yes \\
\hline
\end{tabular}
\end{table}

\subsection{Collaborative Sharing Platforms}

The most prominent collaborative platforms --- MISP and OpenCTI --- inspired
the data model of the present project. MISP organises shared knowledge into
events grouping IOCs, attributes and tags. OpenCTI takes a different
approach by representing the same knowledge as a graph of entities (threat
actors, intrusion sets, malware, indicators) connected by semantic
relationships. Both platforms acknowledge the importance of TLP and of a
strong authentication of contributors, but neither of them offers a
cryptographic proof that the published record has not been altered ---
precisely the gap that the present project aims to fill.

% ---------------- Section 2
\section{Blockchain Technology and Smart Contracts}

\subsection{Definition and Properties}

A blockchain is a distributed data structure in which transactions are
organised in cryptographically linked blocks, replicated across a network
of nodes that share a common consensus protocol. The seminal idea was
introduced by the Bitcoin protocol in 2008 and has since been generalised
to many other applications. Three properties make blockchain particularly
attractive for an integrity-oriented use case:

\begin{itemize}
\item \textbf{Immutability.} Once a block is accepted by the network, any
attempt to modify a past entry would invalidate the cryptographic chain
and would be immediately detected.
\item \textbf{Decentralisation.} The state of the ledger is not under the
control of a single party; multiple nodes hold an authoritative copy.
\item \textbf{Auditability.} Any external party can verify that a given
record has been registered at a given time, without trusting the operator
of the platform.
\end{itemize}

\subsection{Public versus Permissioned Networks}

Two families of blockchain networks coexist. \emph{Public} networks such
as Ethereum are open to any participant, secured by economic incentives and
suitable for use cases where trust between participants is minimal. They
typically charge transaction fees in a native cryptocurrency. \emph{Permissioned}
networks such as Hyperledger Fabric restrict participation to a closed set
of known organisations and rely on identity-based consensus. They are
better suited to enterprise use cases where confidentiality of participation
matters. Table \ref{tab:bc-compare} summarises the comparison.

\begin{table}[!h]
\caption{Comparison of public and permissioned blockchain networks.}
\label{tab:bc-compare}
\centering
\renewcommand{\arraystretch}{1.4}
\begin{tabular}{| C{.20\textwidth} | C{.30\textwidth} | C{.30\textwidth} |}
\hline
\textbf{Criterion} & \textbf{Public (Ethereum)} & \textbf{Permissioned (Fabric)} \\
\hline\hline
Participation     & Open               & Restricted        \\
\hline
Consensus         & Proof of stake     & Identity-based    \\
\hline
Transaction cost  & Paid in ETH        & None              \\
\hline
Throughput        & Moderate           & Higher            \\
\hline
Confidentiality   & Limited            & High              \\
\hline
\end{tabular}
\end{table}

For an academic project, the choice retained is a hybrid one: \textbf{Ganache},
a local Ethereum-compatible network, is used during development for its
accessibility and the maturity of its tooling; \textbf{Hyperledger Fabric}
is identified as the production target because of its alignment with the
collaborative-but-vetted nature of the contributor community.

\subsection{Smart Contracts}

A smart contract is a piece of code deployed on a blockchain that exposes
functions callable through transactions. The execution of the function is
performed by every node of the network, which guarantees that the result is
deterministic and verifiable. In the Ethereum ecosystem, the canonical
language for writing smart contracts is \textbf{Solidity}, a statically
typed language compiled to bytecode for the Ethereum Virtual Machine (EVM).
A smart contract typically exposes:

\begin{itemize}
\item \textbf{Storage variables}, persisted in the state of the contract;
\item \textbf{Functions}, either modifying the state (paid transactions) or
read-only (free calls);
\item \textbf{Events}, emitted during execution and indexable from
off-chain code.
\end{itemize}

In the project, a minimal Solidity contract called \textit{ThreatRegistry}
exposes two write functions (registration and revocation of an IOC) and
emits two events (\texttt{IOCRegistered} and \texttt{IOCRevoked}) that
together form an append-only history --- a pattern known as \emph{event
sourcing}.

\subsection{Off-Chain Integration with Web3}

Most real systems do not store everything on-chain. The recommended pattern
is to keep the raw data in a classical database and to anchor on-chain only
a cryptographic digest of the data. Verification then consists in
re-computing the digest from the off-chain record and comparing it with the
one stored on-chain. The Python library \textbf{Web3.py} is the de-facto
bridge between a backend application and an Ethereum-compatible network: it
allows the backend to read the state of the contract, to build and sign
transactions and to wait for their confirmation.

% ---------------- Section 3
\section{Large Language Models for User Assistance}

\subsection{From Statistical NLP to Large Language Models}

For decades, natural language processing relied on hand-crafted rules and
statistical models trained on relatively small corpora. The introduction of
the \emph{Transformer} architecture in 2017 marked a turning point by
making it possible to train very large models on very large corpora,
giving rise to the family of large language models (LLMs). These models
have demonstrated remarkable abilities at generating coherent text,
explaining concepts in natural language and engaging in multi-turn dialogue.

\subsection{Usage Modes}

LLMs can be used in different modes depending on the level of autonomy
expected from the model, as summarised in Table \ref{tab:llm-modes}.

\begin{table}[!h]
\caption{Common usage modes of large language models.}
\label{tab:llm-modes}
\centering
\renewcommand{\arraystretch}{1.4}
\begin{tabular}{| C{.20\textwidth} | p{.65\textwidth} |}
\hline
\textbf{Mode} & \textbf{Description} \\
\hline\hline
Direct generation & \nohyphens{The model produces an answer from the prompt alone.} \\
\hline
Retrieval-augmented & \nohyphens{A retrieval step extracts a relevant context from a knowledge base; the model is then asked to answer using this context.} \\
\hline
Agent-based & \nohyphens{The model can call external tools (APIs, calculators, search engines) to gather information before answering.} \\
\hline
\end{tabular}
\end{table}

In this project, the assistant follows the \emph{retrieval-augmented} mode
in a controlled way: a small, TLP-filtered context is extracted from the
relational database and injected into the prompt before the call to the
model. This approach preserves the confidentiality guarantees of the
platform while keeping the conversational flexibility of the model.

\subsection{Risks and Mitigations}

Three well-known risks of LLM-based assistants are particularly relevant in
a cybersecurity context: \emph{hallucinations} (the model invents facts),
\emph{prompt injection} (a malicious user tricks the model into ignoring
its system prompt) and \emph{data leakage} (the model echoes sensitive data
present in its context). The mitigations adopted in the project are
respectively to limit the assistant to a basic conceptual scope, to
sanitise the user input on the server side, and to filter the context
through the TLP rules before any call to the model.

\section*{Conclusion}
This second chapter introduced the conceptual foundations on which the
proposed solution is built. Section 1 clarified the discipline of cyber
threat intelligence and the role of indicators of compromise. Section 2
presented blockchain technology and the smart-contract model that will be
used as an integrity anchor. Section 3 outlined how large language models
can be safely embedded in a sensitive application. The next chapter brings
these three pillars together in the concrete design and implementation of
the platform, and evaluates the resulting prototype.

% =============================================================
% CHAPTER 3 — DESIGN, IMPLEMENTATION AND EVALUATION
% =============================================================
\chapter{Solution Design, Implementation and Evaluation}
\label{ch:solution}
\section*{Introduction}
This third and final chapter brings together the project scope defined in
Chapter \ref{ch:scope} and the theoretical foundations introduced in
Chapter \ref{ch:background}. Section 1 presents the design of the system:
the layered architecture, the use-case and class diagrams, the relational
data model and the main interaction sequences. Section 2 details the
implementation choices, the technologies retained and the way each
functional module was realised. Section 3 is devoted to the evaluation of
the delivered prototype, both against the initial objectives and through a
manual test campaign.

% ---------------- Section 1 — Design
\section{Design}

\subsection{Layered Architecture}

The platform follows a strict layered architecture, illustrated in
Figure \ref{fig:archi}. Each layer has a single responsibility and
communicates only with its immediate neighbours.

\begin{figure}[!h]
\centering
\includegraphics[width=0.8\textwidth]{images/architecture.png}
\caption{Four-layer architecture of the platform. Source: authors.}
\label{fig:archi}
\end{figure}

\subsubsection{Presentation Layer}
A single-page application written in React and enriched with the Recharts
library for the interactive visualisations of the public dashboard.

\subsubsection{Application Layer}
A Python backend implemented with the FastAPI framework. It hosts the
business logic, the validation rules, the TLP filtering and the
authentication middleware.

\subsubsection{Conversational Layer}
A thin wrapper around a large language model. It enriches each user
question with a TLP-filtered context extracted from the database before
calling the model.

\subsubsection{Blockchain Layer}
A Ganache local network running an Ethereum-compatible node, hosting the
\textit{ThreatRegistry} smart contract responsible for anchoring validated
indicators.

\subsection{Use-Case Diagram}

Figure \ref{fig:actors-bis} recalls the use-case diagram introduced in
Chapter \ref{ch:scope}. The three categories of users interact with the
platform through clearly differentiated entry points.

\begin{figure}[!h]
\centering
\includegraphics[width=0.85\textwidth]{images/usecase.png}
\caption{Use-case diagram showing the actors and their interactions with the platform.}
\label{fig:actors-bis}
\end{figure}

\subsection{Data Model}

The relational schema contains five core tables, illustrated by the entity
relationship diagram of Figure \ref{fig:erd}.

\begin{figure}[!h]
\centering
\includegraphics[width=0.9\textwidth]{images/erd.png}
\caption{Simplified entity-relationship diagram of the platform.}
\label{fig:erd}
\end{figure}

The central \texttt{iocs} table stores the technical indicators with their
contextual metadata. The \texttt{organisations} table stores the
contributing entities together with a trust score that drives the automatic
validation policy. The \texttt{threat\_actors} and \texttt{malware\_samples}
tables follow the same modelling principles but are dedicated to
higher-level artefacts. Finally, the \texttt{blockchain\_records} table
bridges the relational world and the ledger by storing for each anchored
IOC the transaction hash and the block number returned by the smart
contract.

\subsection{Life-Cycle of an IOC}

The status field of an IOC tracks its life-cycle through the six states
described in Table \ref{tab:status}.

\begin{table}[!h]
\caption{Life-cycle of an IOC and the meaning of each status value.}
\label{tab:status}
\centering
\renewcommand{\arraystretch}{1.4}
\begin{tabular}{| C{.20\textwidth} | p{.70\textwidth} |}
\hline
\textbf{Status} & \textbf{Meaning} \\
\hline\hline
\texttt{validated}      & \nohyphens{Default state after a successful submission; immediately anchored on the blockchain.} \\
\hline
\texttt{pending}        & \nohyphens{Temporary state for low-trust submissions; requires administrator review.} \\
\hline
\texttt{false\_positive} & \nohyphens{Flagged by its own contributor as erroneous; a revocation event is anchored on the chain.} \\
\hline
\texttt{revoked}        & \nohyphens{Definitively invalidated; revocation is logged on the chain.} \\
\hline
\texttt{rejected}       & \nohyphens{Refused for abuse or inappropriate content; not anchored.} \\
\hline
\texttt{deprecated}     & \nohyphens{Kept for historical reference only.} \\
\hline
\end{tabular}
\end{table}

\subsection{Main Interaction Sequence}

Figure \ref{fig:seq} describes the sequence of interactions triggered by the
submission of an IOC: the backend validates the API key, applies the
automatic validation rules, persists the indicator in PostgreSQL, then
calls the smart contract through Web3.py and finally stores the resulting
transaction hash.

\begin{figure}[!h]
\centering
\includegraphics[width=0.95\textwidth]{images/sequence.png}
\caption{Sequence diagram of an IOC submission.}
\label{fig:seq}
\end{figure}

The corresponding decision logic is formalised in Algorithm \ref{algo:submit}.

\begin{algorithm}[!h]
\setstretch{1.15}
\caption{Processing of an incoming IOC submission}
\label{algo:submit}
\SetKwInOut{Input}{Input}
\SetKwInOut{Output}{Output}
\Input{$P$ : the JSON payload of the submission\\
$K$ : the API key carried by the HTTP request}
\Output{$R$ : the persisted IOC together with its blockchain proof}

$\textit{org} \gets \textsf{lookupOrganisation}(K)$\;
\If{$\textit{org} = \emptyset$ \textbf{or} $\textit{org}.\textit{status} \neq$ \textsf{approved}}{
  \textbf{reject} with HTTP 401\;
}
\textsf{validateSchema}($P$)\;
\eIf{$\textit{org}.\textit{trust\_score} \geq 50$ \textbf{and not} \textsf{isSuspicious}($P$)}{
  $\textit{status} \gets$ \texttt{validated}\;
}{
  $\textit{status} \gets$ \texttt{pending}\;
}
$R \gets \textsf{persistIOC}(P, \textit{org}, \textit{status})$\;
\If{$\textit{status} =$ \texttt{validated}}{
  $\textit{tx} \gets \textsf{anchorOnBlockchain}(R)$\;
  $\textsf{persistBlockchainRecord}(R, \textit{tx})$\;
}
\textbf{return} $R$\;
\end{algorithm}

% ---------------- Section 2 — Implementation
\section{Implementation}

\subsection{Technology Stack}

The technologies retained for the project are listed in
Table \ref{tab:stack}. They were selected for their maturity, their
alignment with the requirements and their relevance to the academic
context.

\begin{table}[!h]
\caption{Technology stack of the platform with justification.}
\label{tab:stack}
\centering
\renewcommand{\arraystretch}{1.4}
\begin{tabular}{| C{.20\textwidth} | C{.25\textwidth} | p{.40\textwidth} |}
\hline
\textbf{Component} & \textbf{Technology} & \textbf{Justification} \\
\hline\hline
Frontend         & React + Recharts        & \nohyphens{Mature ecosystem, interactive charts} \\
\hline
Backend          & Python 3.11 + FastAPI   & \nohyphens{Native async, automatic OpenAPI documentation} \\
\hline
Database         & PostgreSQL 15           & \nohyphens{Strong relational guarantees} \\
\hline
ORM              & SQLAlchemy + Alembic    & \nohyphens{Typed queries, clean migrations} \\
\hline
LLM              & Anthropic Claude API    & \nohyphens{High-quality conversational responses} \\
\hline
Enrichment       & ip-api, AbuseIPDB       & \nohyphens{Free IP reputation services} \\
\hline
Blockchain (dev) & Ganache + Truffle       & \nohyphens{Local Ethereum-compatible network} \\
\hline
Blockchain (prod)& Hyperledger Fabric      & \nohyphens{Permissioned ledger for organisations} \\
\hline
Smart contract   & Solidity 0.8            & \nohyphens{Standard EVM language} \\
\hline
Blockchain SDK   & Web3.py 6               & \nohyphens{Python bridge to Ethereum} \\
\hline
Authentication   & JWT (python-jose)       & \nohyphens{Stateless administrator sessions} \\
\hline
\end{tabular}
\end{table}

\subsection{Submission and Collaboration Module}

\subsubsection{Authentication by API Key}
Every contributor request carries an API key in a custom HTTP header
\texttt{X-API-Key}. The key is generated by the administrator at the moment
of approval, transmitted once by e-mail and stored hashed (SHA-256) in the
database. The backend rejects any request whose hash does not match a known
and active organisation. This mechanism, summarised in
Table \ref{tab:apikey}, intentionally trades the comfort of a classical
session for the precision of per-request traceability.

\begin{table}[!h]
\caption{Properties of the API key authentication mechanism.}
\label{tab:apikey}
\centering
\renewcommand{\arraystretch}{1.4}
\begin{tabular}{| C{.20\textwidth} | p{.70\textwidth} |}
\hline
\textbf{Property} & \textbf{Implementation} \\
\hline\hline
Format        & \nohyphens{SHA-256 hexadecimal string} \\
\hline
Transmission  & \nohyphens{Custom HTTP header X-API-Key} \\
\hline
Storage       & \nohyphens{Hashed in the organisations table, never in clear} \\
\hline
Revocation    & \nohyphens{Instant by changing the status of the organisation} \\
\hline
Renewal       & \nohyphens{New key generated by the administrator and sent by e-mail} \\
\hline
\end{tabular}
\end{table}

\subsubsection{Submission Endpoints}
Three structured endpoints expose the submission of, respectively,
indicators, threat actors and malware samples. Their JSON payloads carry
the same conceptual structure: identification of the artefact, contextual
description, TLP level and optional enrichment data. The backend validates
each payload through a Pydantic schema before persistence.

\subsubsection{Automatic Validation and the Trust Score}
To preserve scalability, the platform does not rely on manual validation of
every submission. Each organisation carries a numerical trust score between
0 and 100, computed from its history. A submission whose trust score is
above the threshold is automatically validated and immediately anchored. A
submission below the threshold or matching a suspicious pattern is placed
in a pending state for administrator review.

\subsubsection{False-Positive Workflow}
A dedicated endpoint allows a contributor to flag one of its own indicators
as a false positive. The backend verifies ownership, updates the status of
the IOC to \texttt{false\_positive} and emits a revocation event on the
blockchain.

\subsection{Conversational AI Assistance Module}

\subsubsection{Scope of the Assistant}
The chatbot integrated into the platform is intentionally limited to an
assistance role. In the current prototype it relies on a local configuration
so as to ensure its availability independently from any external network
condition. Its core mission is to answer basic, factual questions about the
cybersecurity vocabulary used on the platform --- for example explaining
what the Traffic Light Protocol is, what is meant by an indicator of
compromise or how to read a blockchain proof. It does \emph{not} compute
any danger score, does not participate in the validation of indicators and
never reveals data classified above the visitor's permission level.

\subsubsection{Architecture of the Chatbot}
The assistant follows the four-step pipeline shown in
Figure \ref{fig:chatbot}.

\begin{figure}[!h]
\centering
\includegraphics[width=0.95\textwidth]{images/chatbot.png}
\caption{Processing pipeline of a user question by the conversational assistant.}
\label{fig:chatbot}
\end{figure}

\subsubsection{Contextual Filtering}
A fundamental property of the implementation is that the assistant never
has direct access to the database. The backend extracts a tailored context
based on the question and filters it according to the user's permission
level before sending it to the language model. This indirection enforces
TLP rules even on the AI layer.

\subsubsection{Implemented Capabilities}
The capabilities currently delivered are summarised in Table \ref{tab:chat}.

\begin{table}[!h]
\caption{Capabilities currently implemented in the conversational assistant.}
\label{tab:chat}
\centering
\renewcommand{\arraystretch}{1.4}
\begin{tabular}{| C{.25\textwidth} | p{.65\textwidth} |}
\hline
\textbf{Capability} & \textbf{Description} \\
\hline\hline
Concept explanation & \nohyphens{Plain-language definition of cybersecurity terms (IOC, TLP, APT, etc.).} \\
\hline
Database guidance & \nohyphens{High-level orientation about where to find a piece of information on the platform.} \\
\hline
Reading assistance & \nohyphens{Help in interpreting the fields of an indicator or the meaning of a blockchain proof.} \\
\hline
Multi-turn memory & \nohyphens{Short conversational history kept inside the current session.} \\
\hline
\end{tabular}
\end{table}

\subsection{Blockchain Anchoring Module}

\subsubsection{Scope of the Integration}
In the current prototype, blockchain anchoring is implemented \textbf{only
for indicators of compromise}. Threat actor profiles and malware samples
are stored in the database and exposed by the API, but they are not yet
pushed to the ledger. Extending the anchoring mechanism to these two
additional categories is identified as a perspective.

\subsubsection{The ThreatRegistry Smart Contract}
The smart contract, named \textit{ThreatRegistry}, is written in Solidity
0.8 and deployed on a local Ganache network. Its responsibility is
intentionally narrow: it persists a minimal record per IOC and emits the
events needed to reconstruct the full history. The fields anchored on the
chain are summarised in Table \ref{tab:onchain}.

\begin{table}[!h]
\caption{Data fields recorded on the blockchain for each IOC.}
\label{tab:onchain}
\centering
\renewcommand{\arraystretch}{1.4}
\begin{tabular}{| C{.18\textwidth} | C{.12\textwidth} | p{.55\textwidth} |}
\hline
\textbf{Field} & \textbf{Type} & \textbf{Role} \\
\hline\hline
\texttt{iocHash}   & string  & \nohyphens{SHA-256 of the raw value (never the value itself)} \\
\hline
\texttt{category}  & string  & \nohyphens{Category of threat (phishing, botnet, scanner, etc.)} \\
\hline
\texttt{tlp}       & string  & \nohyphens{TLP level of the indicator} \\
\hline
\texttt{timestamp} & uint256 & \nohyphens{Block timestamp, used as a tamper-evident date} \\
\hline
\end{tabular}
\end{table}

A deliberate design choice is to anchor only the hash of the indicator and
not its raw value. This decision protects the confidentiality of the
underlying data while preserving the ability of any external party to
verify that a value presented to them matches the one originally registered.

\subsubsection{Backend--Blockchain Integration}
The backend communicates with the smart contract through the Web3.py
library. When the validation step decides that an IOC is acceptable, the
backend builds the registration transaction, sends it, waits for the
receipt and finally stores the resulting transaction hash and block number
in the \texttt{blockchain\_records} table. The combination of the database
record and the on-chain entry forms the proof of integrity that is later
displayed on the public detail page.

\subsubsection{Event Sourcing and Immutability}
Once an entry has been written on the chain it cannot be modified or
deleted. Correcting an erroneous indicator therefore relies on an
event-sourcing model: a registration transaction is followed, when needed,
by a revocation transaction referring to the same identifier. The frontend
reconstructs the current state of the indicator by replaying the sequence
of events.

\subsection{User Interfaces}

\subsubsection{Public Dashboard}
The public dashboard is the entry point for any visitor. It displays a
paginated list of indicators classified as TLP:GREEN or TLP:WHITE, a
search bar, two charts dedicated respectively to the distribution by
category and to the temporal evolution of submissions, an access point to
the AI assistant and a link to the contributor registration form.

\subsubsection{Indicator Detail Page}
The detail page of an indicator shows the contextual data submitted by the
contributor, the enrichment information automatically retrieved from
external services, a coloured TLP badge, the current status of the
indicator and the full sequence of blockchain events that affected it. The
transaction hash is displayed together with the block number and the
corresponding timestamp.

\subsubsection{Contributor Registration Form}
The registration form collects the legal information needed to verify the
identity of a new organisation: legal name, registration number, technical
contact e-mail, website, activity description and country.

\subsubsection{Administration Dashboard}
The administration dashboard, accessible only after a JWT-based
authentication, exposes the list of pending applications, the actions to
approve or refuse them, the management of API keys and a global view of
the platform statistics.

% ---------------- Section 3 — Evaluation
\section{Evaluation}

\subsection{Coverage of the Initial Objectives}

Table \ref{tab:eval-obj} measures how each general objective declared in
Chapter \ref{ch:scope} has been covered by the delivered prototype.

\begin{table}[!h]
\caption{Evaluation of the prototype against the general objectives.}
\label{tab:eval-obj}
\centering
\renewcommand{\arraystretch}{1.35}
\begin{tabular}{| C{.18\textwidth} | p{.55\textwidth} | C{.12\textwidth} |}
\hline
\textbf{Objective} & \textbf{Outcome} & \textbf{Status} \\
\hline\hline
Centralisation     & \nohyphens{Submission endpoints for IOCs, threat actors and malware samples are fully operational.} & Achieved \\
\hline
Accessibility      & \nohyphens{Public dashboard exposes TLP:GREEN/WHITE data without authentication.} & Achieved \\
\hline
AI assistance      & \nohyphens{Local chatbot answers basic questions on the platform vocabulary.} & Achieved \\
\hline
Integrity          & \nohyphens{IOCs are anchored on a Ganache network through the ThreatRegistry contract.} & Achieved \\
\hline
Security           & \nohyphens{API key authentication with hashed storage and revocation is in place.} & Achieved \\
\hline
Compliance         & \nohyphens{TLP filtering is enforced on every endpoint, including the chatbot context.} & Achieved \\
\hline
Full anchoring     & \nohyphens{Threat actors and malware samples are not yet anchored on the chain.} & Partial \\
\hline
\end{tabular}
\end{table}

\subsection{Coverage of the Functional Requirements}

The ten functional requirements introduced in Chapter \ref{ch:scope} are
revisited in Table \ref{tab:eval-fr}.

\begin{table}[!h]
\caption{Coverage of the functional requirements by the prototype.}
\label{tab:eval-fr}
\centering
\renewcommand{\arraystretch}{1.35}
\begin{tabular}{| C{.08\textwidth} | p{.65\textwidth} | C{.17\textwidth} |}
\hline
\textbf{Code} & \textbf{Requirement} & \textbf{Status} \\
\hline\hline
FR-1  & \nohyphens{Submit an IOC}                          & Done \\
\hline
FR-2  & \nohyphens{Submit a threat actor}                  & Done (no anchoring) \\
\hline
FR-3  & \nohyphens{Submit a malware sample}                & Done (no anchoring) \\
\hline
FR-4  & \nohyphens{Browse and search the catalogue}        & Done \\
\hline
FR-5  & \nohyphens{Consult the detail of an indicator}     & Done \\
\hline
FR-6  & \nohyphens{Interact with the AI assistant}         & Done (basic scope) \\
\hline
FR-7  & \nohyphens{Verify the blockchain proof}            & Done \\
\hline
FR-8  & \nohyphens{Flag a false positive}                  & Done \\
\hline
FR-9  & \nohyphens{Approve or revoke a contributor}        & Done \\
\hline
FR-10 & \nohyphens{Review pending submissions}             & Done \\
\hline
\end{tabular}
\end{table}

\subsection{Functional Validation Campaign}

Functional validation was carried out through a manual test plan covering
the four user flows --- public consultation, contributor submission,
false-positive flagging and administrator approval. The results are
summarised in Table \ref{tab:tests}.

\begin{table}[!h]
\caption{Summary of the functional validation campaign.}
\label{tab:tests}
\centering
\renewcommand{\arraystretch}{1.35}
\begin{tabular}{| C{.08\textwidth} | p{.65\textwidth} | C{.17\textwidth} |}
\hline
\textbf{ID} & \textbf{Scenario} & \textbf{Result} \\
\hline\hline
T1 & \nohyphens{Anonymous visitor consults the public dashboard}          & Pass \\
\hline
T2 & \nohyphens{Anonymous visitor uses the chatbot}                       & Pass \\
\hline
T3 & \nohyphens{Contributor submits a valid IOC}                          & Pass \\
\hline
T4 & \nohyphens{Contributor submits with an invalid API key}              & Reject (expected) \\
\hline
T5 & \nohyphens{Contributor flags one of its own IOCs as false positive}  & Pass \\
\hline
T6 & \nohyphens{Contributor tries to flag an IOC of another organisation} & Reject (expected) \\
\hline
T7 & \nohyphens{Administrator approves a new contributor application}     & Pass \\
\hline
T8 & \nohyphens{Visitor verifies the blockchain proof of an indicator}    & Pass \\
\hline
\end{tabular}
\end{table}

\subsection{Strengths and Limitations}

The evaluation highlights three particular strengths of the delivered work.
First, the architecture is clean and modular: each layer can evolve
independently. Second, the blockchain integration is functional end-to-end,
which is a non-trivial achievement at academic level given the learning
curve of Solidity and Web3. Third, the chatbot demonstrates how an AI
layer can be added to a sensitive application without breaking its security
model, thanks to the systematic filtering of the context before any call to
the language model.

The team also identifies a number of limitations. The blockchain anchoring
is currently restricted to IOCs; the chatbot covers only basic
concept-explanation questions; and the prototype has not yet been
benchmarked under a heavy concurrent load. These limitations are not
disqualifying for the academic exercise but constitute clear leads for
continuation, discussed in the general conclusion.

\section*{Conclusion}
This third chapter walked through the design, the implementation and the
evaluation of the platform. The design section translated the requirements
of Chapter \ref{ch:scope} and the theoretical pillars of
Chapter \ref{ch:background} into a coherent architecture, a relational
data model and a clear life-cycle for indicators. The implementation section
detailed the technology stack and the way each functional module was
realised. The evaluation section measured the prototype against the initial
objectives and confirmed that the principal targets have been reached, with
the documented exception of the full anchoring of all artefact categories.

% =============================================================
% GENERAL CONCLUSION
% =============================================================
\chapter*{General Conclusion and Perspectives}
\addcontentsline{toc}{chapter}{General Conclusion and Perspectives}
\markboth{General Conclusion and Perspectives}{}

This Personal Professional Project led us, in the space of fourteen weeks,
to design and implement a complete collaborative threat intelligence
platform combining a structured submission workflow, a conversational AI
assistant and an integrity layer based on blockchain technology. Beyond the
technical realisation, the project gave us the opportunity to address a
real cybersecurity challenge --- the trustworthy sharing of indicators of
compromise --- and to confront a number of decisions that have practical
implications.

From a technical standpoint, the project allowed us to articulate four
heterogeneous worlds: a modern web frontend in React, an asynchronous
Python backend with FastAPI, an Ethereum-compatible blockchain network
with Solidity and Web3, and a large language model accessed through a
dedicated API. The discipline imposed by the layered architecture and by
the careful enforcement of the Traffic Light Protocol across every layer
was probably the most valuable lesson of the project: security is not a
feature added at the end but a property that emerges from consistent design
choices.

From a methodological standpoint, the structured organisation of the work
into five phases and the weekly synchronisation meetings allowed the four
members of the team to converge despite the diversity of the topics handled
in parallel. We also learned the importance of an honest evaluation:
rather than overstating the scope of the prototype, we explicitly delimit
what is functional today --- IOC anchoring, basic chatbot, full
administrative workflow --- and what remains to be done.

\paragraph{Perspectives.} Several extensions are envisaged to continue the
project beyond its initial scope:

\begin{itemize}
\item \textbf{Extending the blockchain anchoring.} The current prototype
anchors only indicators of compromise. A natural and high-value extension
consists in applying the same anchoring mechanism to threat actor profiles
and to malware samples, so that the integrity guarantees provided by the
ledger cover the entire knowledge base.
\item \textbf{Enriching the conversational assistant.} The chatbot is
deliberately limited today to basic explanations such as the meaning of
the Traffic Light Protocol. Future iterations could allow it to summarise
complex indicators, to generate targeted mitigation advice or to support
multi-language interaction, while preserving the strict TLP filtering on
the context provided to the model.
\item \textbf{Migration to a permissioned ledger.} The development was
conducted on Ganache for reasons of accessibility. A migration to
Hyperledger Fabric, anticipated in the design, would bring the project
closer to a deployable production system, with explicit identity management
of participating organisations.
\item \textbf{Performance benchmarking.} The platform has not yet been
evaluated under heavy concurrent load. A systematic benchmarking campaign
would allow us to characterise the practical limits of the chosen
architecture.
\item \textbf{Automated trust scoring.} The trust score that drives the
automatic validation policy is, in the current prototype, configured by
the administrator. A logical evolution is to compute it dynamically from
the historical behaviour of each contributor.
\item \textbf{Integration with public threat feeds.} Finally, the platform
could be enriched by an ingestion of open feeds such as MISP, Feodo
Tracker or AbuseIPDB, which would broaden the catalogue available to
visitors.
\end{itemize}

The project, while academic in scope, demonstrates that the combination of
artificial intelligence and blockchain technology can deliver tangible
value in the field of collaborative cybersecurity. It also reinforced our
conviction that the design of secure systems is, above all, a matter of
careful coordination between functional needs, technical choices and human
practices.

% =============================================================
% APPENDIX
% =============================================================
\begin{appendix}
\chapter{Summary of Public API Endpoints}

Table \ref{tab:endpoints} recalls the main endpoints exposed by the
platform. The full OpenAPI specification is automatically generated by
FastAPI and accessible through the \texttt{/docs} route at runtime.

\begin{table}[!h]
\caption{Main HTTP endpoints exposed by the platform.}
\label{tab:endpoints}
\centering
\renewcommand{\arraystretch}{1.35}
\begin{tabular}{| C{.08\textwidth} | C{.30\textwidth} | p{.40\textwidth} | C{.12\textwidth} |}
\hline
\textbf{Method} & \textbf{Route} & \textbf{Purpose} & \textbf{Auth} \\
\hline\hline
GET  & /iocs                  & \nohyphens{List public IOCs}                  & Public \\
\hline
GET  & /iocs/:id              & \nohyphens{Detail of an IOC}                  & Public \\
\hline
POST & /iocs/submit           & \nohyphens{Submit a new IOC}                  & API key \\
\hline
GET  & /iocs/mine             & \nohyphens{List own submissions}              & API key \\
\hline
POST & /iocs/false-positive   & \nohyphens{Flag an IOC as false positive}     & API key \\
\hline
GET  & /threat-actors         & \nohyphens{List public threat actors}         & Public \\
\hline
POST & /threat-actors/submit  & \nohyphens{Submit a threat actor}             & API key \\
\hline
GET  & /malware               & \nohyphens{List public malware samples}       & Public \\
\hline
POST & /malware/submit        & \nohyphens{Submit a malware sample}           & API key \\
\hline
POST & /register              & \nohyphens{Contributor registration request}  & Public \\
\hline
POST & /admin/login           & \nohyphens{Administrator authentication}      & Public \\
\hline
POST & /admin/approve/:id     & \nohyphens{Approve a registration}            & Admin JWT \\
\hline
POST & /admin/revoke/:id      & \nohyphens{Revoke an API key}                 & Admin JWT \\
\hline
GET  & /blockchain/verify/:id & \nohyphens{Verify the blockchain proof}       & Public \\
\hline
POST & /chat                  & \nohyphens{Interact with the AI assistant}    & Public \\
\hline
\end{tabular}
\end{table}

\end{appendix}

% =============================================================
% NETOGRAPHY
% =============================================================
\chapter*{Netography}
\addcontentsline{toc}{chapter}{Netography}
\markboth{Netography}{}

\begin{enumerate}[label={[\arabic*]}]

\item MISP Project, \textit{Open Source Threat Intelligence Platform}.
\url{https://www.misp-project.org/} (accessed May 2026).

\item OpenCTI Platform, \textit{Open Cyber Threat Intelligence Platform}.
\url{https://www.opencti.io/} (accessed May 2026).

\item FIRST.org, \textit{Traffic Light Protocol (TLP) Version 2.0}.
\url{https://www.first.org/tlp/} (accessed May 2026).

\item AbuseIPDB, \textit{IP Address Abuse Database}.
\url{https://www.abuseipdb.com/} (accessed May 2026).

\item Feodo Tracker, \textit{Abuse.ch Botnet C2 IP Blocklist}.
\url{https://feodotracker.abuse.ch/} (accessed May 2026).

\item VirusTotal, \textit{File and URL Analyser}.
\url{https://www.virustotal.com/} (accessed May 2026).

\item FastAPI Documentation, \textit{Modern Python Web Framework for APIs}.
\url{https://fastapi.tiangolo.com/} (accessed May 2026).

\item PostgreSQL Global Development Group, \textit{PostgreSQL Documentation}.
\url{https://www.postgresql.org/docs/} (accessed May 2026).

\item SQLAlchemy, \textit{The Database Toolkit for Python}.
\url{https://www.sqlalchemy.org/} (accessed May 2026).

\item React, \textit{A JavaScript Library for Building User Interfaces}.
\url{https://react.dev/} (accessed May 2026).

\item Recharts, \textit{Charting Library Built with React}.
\url{https://recharts.org/} (accessed May 2026).

\item Solidity Programming Language, \textit{Official Documentation}.
\url{https://soliditylang.org/} (accessed May 2026).

\item Truffle Suite, \textit{Smart Contract Development Tools}.
\url{https://trufflesuite.com/} (accessed May 2026).

\item Ganache, \textit{Personal Ethereum Blockchain for Development}.
\url{https://trufflesuite.com/ganache/} (accessed May 2026).

\item Web3.py, \textit{Python Library for Interacting with Ethereum}.
\url{https://web3py.readthedocs.io/} (accessed May 2026).

\item Hyperledger Fabric, \textit{Enterprise Permissioned Blockchain Framework}.
\url{https://www.hyperledger.org/projects/fabric} (accessed May 2026).

\item Anthropic, \textit{Claude API Documentation}.
\url{https://docs.anthropic.com/} (accessed May 2026).

\item ENISA, \textit{Threat Landscape Report}.
\url{https://www.enisa.europa.eu/} (accessed May 2026).

\item MITRE Corporation, \textit{ATT\&CK Knowledge Base}.
\url{https://attack.mitre.org/} (accessed May 2026).

\item NIST, \textit{Guide to Cyber Threat Information Sharing} (SP 800-150).
\url{https://csrc.nist.gov/publications/detail/sp/800-150/final} (accessed May 2026).

\end{enumerate}

\end{document}
